\section{5~V converter derivation}\label{sec:5v-converter}

This section derives the required characteristics of the 5~V rail and verifies the
populated Rev-A implementation: U401 (AP63205WU, \cref{sec:ap632xx}) with L401
(\SI{4.7}{\micro\henry}), C401 (\SI{10}{\micro\farad} input), C402 (\SI{100}{\nano\farad}
BST), and C403/C404 ($2\times\SI{22}{\micro\farad}$ output). The exported netlist confirms
the \texttt{+5V} net contains exactly: the converter output components (L401, C403, C404),
the four K155ID1 $V_{CC}$ pins (U1--U4, pin 5), their per-package bypass capacitors
C1--C4 (\SI{100}{\nano\farad} each), and the power-indicator branch R401
(\SI{5.1}{\kilo\ohm}) feeding D401 to ground --- nothing else. The rail powers the BCD
decoder/drivers and its own indicator LED alone; it has no connection to the MCU or any
other subsystem. This is a
far lighter design point than the flyback of \cref{sec:hv-converter}: a fixed-output
integrated buck, a static logic load two decades below the converter's rating, and no
external compensation, sensing, or clamp networks to verify. The derivation is
correspondingly shorter; where a margin is simply large, it is reported as large and left
alone.

\subsection{Load budget}\label{sec:5v-load}

Each K155ID1 draws $I_{CC} = 16$~mA typical, \SI{25}{\milli\ampere} maximum
(\cref{sec:bcd-driver}; specified at $V_{CC}=$ max, inputs grounded, outputs open). The
cathode current sunk by the selected output does \emph{not} add to this: the driver output
is an open-collector TTL stage sinking current that arrives from the 170~V rail through
the tube, into the output pin, and out the ground pin --- it never passes through
$V_{CC}$. The rail load is therefore the four packages' logic supply current only:
\begin{align*}
  I_{\mathrm{load,typ}} &= 4 \times \SI{16}{\milli\ampere} = \SI{64}{\milli\ampere}, &
  I_{\mathrm{load,max}} &= 4 \times \SI{25}{\milli\ampere} = \SI{100}{\milli\ampere},
\end{align*}
\begin{align*}
  P_{\mathrm{out,typ}} &= 5 \times 0.064 = \SI{0.32}{\watt}, &
  P_{\mathrm{out,max}} &= 5 \times 0.100 = \SI{0.50}{\watt}.
\end{align*}
Bipolar TTL supply current is set by the internal bias networks, not by output state, so
this load is essentially static: it does not change with which digit is lit, and the only
input transitions are the once-per-second (and slower) BCD updates. There is no transient
load profile to size against --- the DC maximum is the whole story. The
\SI{100}{\nano\farad} per-package bypass capacitors C1--C4 are the standard one-per-TTL-package
practice, generously adequate at these switching rates, and contribute no DC load.

One further branch loads the rail continuously: the power-indicator LED D401, fed
through R401 (\SI{5.1}{\kilo\ohm}). \emph{Assumption:} forward voltage
$\approx\SI{2.0}{\volt}$ (generic 0805 part, color unrecorded; the value field reads
``+5V,'' a functional label rather than a part choice). Then
\begin{align*}
  I_{\mathrm{D401}} &= \frac{5-2.0}{\SI{5.1}{\kilo\ohm}} = \SI{0.59}{\milli\ampere}
    \quad\text{(continuous, non-switching)},
\end{align*}
raising the rail totals to \SI{65}{\milli\ampere} typical / \SI{101}{\milli\ampere}
maximum. The increment is 0.6\,\% of the driver load --- below the rounding of every
downstream figure and well inside the $\eta=0.85$ conservatism --- so the derivations
that follow carry the driver-only \SI{100}{\milli\ampere} bound unchanged.

\subsection{Operating point and conduction mode}\label{sec:5v-op}

The input is $V_{\mathrm{BUS}}$, the USB-PD-negotiated rail (9 or \SI{12}{\volt} nominal
contracts; analysis window \SIrange{8.5}{13}{\volt}), tied directly to both IN and EN with
no series element. The window sits deep inside the AP63205's recommended
\SIrange{3.8}{32}{\volt} operating range (\cref{sec:ap632xx}), with \SI{22}{\volt} to the
\SI{35}{\volt} absolute maximum. The ideal buck duty cycle $D = V_{OUT}/V_{IN}$ spans
\begin{align*}
  D_{\min} &= \frac{5}{13} = 0.385 \quad (V_{IN}=\SI{13}{\volt}), &
  D_{\max} &= \frac{5}{8.5} = 0.588 \quad (V_{IN}=\SI{8.5}{\volt}).
\end{align*}
Conduction-loss corrections to $D$ are of order $I_{\mathrm{load}} R_{DS(\mathrm{on})}
\approx \SI{13}{\milli\volt}$ against a 5~V output --- below 0.3\,\% --- and are ignored
throughout. Because $D_{\max} > 0.5$ at the bottom of the window, slope compensation is
required for a peak-current-mode part; the AP63205 integrates it internally and no
external action exists to take (\ds{AP632XX.pdf}{9}). Dropout is not remotely in play:
minimum headroom is $8.5 - 5 = \SI{3.5}{\volt}$, so the near-100\%-duty ``LDO mode''
described qualitatively in \cref{sec:ap632xx} is never entered on this rail.

\paragraph{Conduction mode.} With the populated $L = \SI{4.7}{\micro\henry}$ at
$f_{SW} = \SI{1.1}{\mega\hertz}$, the inductor ripple current in fixed-frequency
operation would be
\begin{align*}
  \Delta I_L \big|_{\SI{8.5}{\volt}} &=
    \frac{5\,(8.5-5)}{8.5 \times \SI{4.7}{\micro\henry} \times \SI{1.1}{\mega\hertz}}
    = \frac{17.5}{43.9} = \SI{0.398}{\ampere}, \\
  \Delta I_L \big|_{\SI{13}{\volt}} &=
    \frac{5\,(13-5)}{13 \times \SI{4.7}{\micro\henry} \times \SI{1.1}{\mega\hertz}}
    = \frac{40}{67.2} = \SI{0.595}{\ampere},
\end{align*}
so the CCM boundary load ($\Delta I_L/2$) is \SIrange{199}{298}{\milli\ampere} across the
input window. The actual load, \SIrange{64}{100}{\milli\ampere}, sits below the boundary
at every corner: the converter does not run continuous fixed-frequency PWM here. This is
not the design concern it was for the flyback, where DCM was a \emph{designed-in}
requirement that the compensation, current limit, and transformer reset all depended on
(\cref{sec:hv-dcm}). In the AP63205 the light-load behavior is an automatic, internally
managed PFM mode: the COMP node is clamped at a floor corresponding to a
\SI{450}{\milli\ampere} peak inductor current, zero-cross detection on the low-side FET
blocks reverse current, and the part skips cycles as the load falls, maintaining
regulation continuously through the transition (\ds{AP632XX.pdf}{10}). No component
choice in this design selects or depends on the mode; it is reported here as an operating
fact, not verified as a requirement.

The PFM operating point follows directly from the \SI{450}{\milli\ampere} clamp.
Each pulse ramps the inductor to $I_{\mathrm{pk}} = \SI{0.45}{\ampere}$ and back to zero:
\begin{align*}
  t_{\mathrm{on}} &= \frac{L\,I_{\mathrm{pk}}}{V_{IN}-V_{OUT}}
    = \frac{\SI{2.115}{\micro\volt\second}}{\SI{8}{\volt}} = \SI{264}{\nano\second}
    \;\; (V_{IN}=\SI{13}{\volt}); \quad
    \frac{\SI{2.115}{\micro\volt\second}}{\SI{3.5}{\volt}} = \SI{604}{\nano\second}
    \;\; (V_{IN}=\SI{8.5}{\volt}), \\
  t_{\mathrm{dem}} &= \frac{L\,I_{\mathrm{pk}}}{V_{OUT}}
    = \frac{\SI{2.115}{\micro\volt\second}}{\SI{5}{\volt}} = \SI{423}{\nano\second}
    \quad (\text{both corners}),
\end{align*}
delivering charge $Q = \tfrac{1}{2} I_{\mathrm{pk}} (t_{\mathrm{on}}+t_{\mathrm{dem}})
= \SI{155}{\nano\coulomb}$ (\SI{13}{\volt}) to \SI{231}{\nano\coulomb}
(\SI{8.5}{\volt}) per pulse. The average pulse repetition rate needed to support the load
is $f_p = I_{\mathrm{load}}/Q$:
\begin{align*}
  f_p &= \frac{0.100}{155\,\si{\nano\coulomb}} \approx \SI{650}{\kilo\hertz}
    \;\;(\SI{100}{\milli\ampere}, \SI{13}{\volt})
  &&\text{down to}&
  f_p &= \frac{0.064}{231\,\si{\nano\coulomb}} \approx \SI{277}{\kilo\hertz}
    \;\;(\SI{64}{\milli\ampere}, \SI{8.5}{\volt}),
\end{align*}
i.e., roughly a quarter to two-thirds of the nominal clock, well above audibility and an
idealized average (real skipping is bursty). The datasheet's typical-efficiency curve
($V_{IN}=\SI{12}{\volt}$, $V_{OUT}=\SI{5}{\volt}$) reads approximately 88--90\,\% in the
60--100~mA region (\ds{AP632XX.pdf}{2}); $\eta = 0.85$ is carried below as a conservative
figure. Average input current is then at most
\begin{align*}
  I_{IN} &= \frac{P_{\mathrm{out,max}}}{\eta\,V_{IN}}
          = \frac{0.50}{0.85 \times 8.5} = \SI{69}{\milli\ampere},
\end{align*}
a negligible increment on the $V_{\mathrm{BUS}}$ budget next to the flyback's draw.

\subsection{Minimum on-time}\label{sec:5v-ontime}

For a buck the on-time shrinks as $V_{IN}$ rises, so the binding corner is
$V_{IN}=\SI{13}{\volt}$ --- the opposite corner from the flyback's. In hypothetical
fixed-frequency operation at the +6\,\% spread-spectrum extreme
($f_{SW} = \SI{1.166}{\mega\hertz}$):
\begin{align*}
  t_{\mathrm{on}} &= \frac{D_{\min}}{f_{SW}}
    = \frac{0.385}{\SI{1.166}{\mega\hertz}} = \SI{330}{\nano\second}
    \quad\text{vs.}\quad t_{\mathrm{on,min}} = \SI{80}{\nano\second}
    \;\;(\cref{sec:ap632xx}),
\end{align*}
a $4.1\times$ margin. In the PFM mode actually occupied, the shortest pulse is the
\SI{264}{\nano\second} on-time computed above --- still $3.3\times$ the minimum. For the
constraint ever to bind at this output voltage, the input would need
$D = t_{\mathrm{on,min}} f_{SW} = 0.093$, i.e.\ $V_{IN} = 5/0.093 \approx \SI{54}{\volt}$,
beyond even the \SI{35}{\volt} absolute maximum. Minimum on-time cannot constrain a 5~V
output anywhere in this part's operating space. Match, no further analysis warranted.

\subsection{Inductor}\label{sec:5v-inductor}

The datasheet's selection equation (\cref{sec:ap632xx}) with its prescribed ripple target
--- $\Delta I_L$ equal to 30--50\,\% of the \emph{2~A rated} load, i.e.\
\SIrange{0.6}{1.0}{\ampere}; the reference is the part's rating, not this application's
100~mA --- gives, at the maximum-ripple corner $V_{IN}=\SI{13}{\volt}$:
\begin{align*}
  L\big|_{\Delta I_L = \SI{1.0}{\ampere}} &=
    \frac{5\,(13-5)}{13 \times 1.0 \times \SI{1.1}{\mega\hertz}} = \SI{2.80}{\micro\henry},
  &
  L\big|_{\Delta I_L = \SI{0.6}{\ampere}} &=
    \frac{40}{13 \times 0.6 \times \SI{1.1}{\mega\hertz}} = \SI{4.66}{\micro\henry},
\end{align*}
and \SIrange{1.87}{3.12}{\micro\henry} at \SI{8.5}{\volt}. The populated
\SI{4.7}{\micro\henry} sits at the large-inductance edge of the 13~V window (ripple
29.8\,\% of 2~A at \SI{13}{\volt}, 19.9\,\% at \SI{8.5}{\volt} --- at or slightly below
the band's 30\,\% floor, the benign direction, which the datasheet itself endorses:
``use a larger inductance for improved efficiency under light load,'' precisely this
rail's condition). More directly: \SI{4.7}{\micro\henry} is the datasheet's own Table-3
BOM value for the AP63205 5~V output (\cref{sec:ap632xx}), and the surrounding parts ---
\SI{10}{\micro\farad} input, $2\times\SI{22}{\micro\farad}$ output,
\SI{100}{\nano\farad} BST --- reproduce that table line for line. The populated design is
the manufacturer's reference BOM applied verbatim, not an independent coincidence.
It also sits inside the general \SIrange{2.2}{10}{\micro\henry} recommended range. Match.

\paragraph{Saturation current.} The schematic records only the footprint
(\texttt{L\_APV\_ANR6045}, a $6.0\times6.0\times\SI{4.5}{\milli\metre}$ shielded power
inductor), not a part number, so the rating is checked against the class, with the
assumption stated: \emph{assume} $I_{\mathrm{sat}} \geq \SI{3.0}{\ampere}$ and DCR
$\leq \SI{50}{\milli\ohm}$, typical of \SI{4.7}{\micro\henry} parts in this size class
(commodity 6045-package \SI{4.7}{\micro\henry} shielded inductors commonly rate
3.5--5~A saturation, 20--50~m$\Omega$ DCR). Against the datasheet's application rule
($\geq$35\,\% above maximum load: $1.35 \times 0.100 = \SI{135}{\milli\ampere}$) the
assumed rating is $>20\times$; against the largest current the inductor ever carries in
normal operation --- the \SI{450}{\milli\ampere} PFM peak clamp --- it is $\approx 7\times$;
and it covers even the worst-case fault current, the \SI{3.1}{\ampere} maximum HS peak
limit, at or near its hard-saturation rating, with the converter's hiccup protection
(current-limit events for 2~ms force a 16~ms shutdown, \ds{AP632XX.pdf}{11}) bounding any
fault dwell time. The DCR assumption satisfies the $<\SI{100}{\milli\ohm}$ efficiency
guidance. Match under the stated assumption; the actual part's saturation rating should
be confirmed when the BOM is finalized, the same status as the unrecorded capacitor
series in \cref{sec:hv-caps}.

\subsection{Current-limit margin}\label{sec:5v-ilim}

The operative cycle-by-cycle limit is the HS peak current limit, \SI{2.5}{\ampere}
minimum (\cref{sec:ap632xx}, which also records why the higher ``LS valley'' figure is
not used). The largest peak inductor current the control loop commands at this load is
the \SI{450}{\milli\ampere} PFM clamp:
\begin{align*}
  \frac{I_{\mathrm{lim}}^{\min}}{I_{\mathrm{pk}}} &= \frac{2.5}{0.45} = 5.6\times.
\end{align*}
Even if the load quadrupled and the converter re-entered CCM, the peak
($I_{\mathrm{load}} + \Delta I_L/2 \approx \SI{0.7}{\ampere}$ at \SI{400}{\milli\ampere}
load) would still leave $>3.5\times$. The current limit exists on this rail purely as
fault protection --- a shorted +5~V net lands in hiccup mode rather than anything
thermal --- and no operating condition approaches it. Large margin; nothing further to
derive.

\subsection{Output capacitance and rail accuracy}\label{sec:5v-cout}

The populated $2\times\SI{22}{\micro\farad}$ (C403/C404) equals the datasheet BOM example
and its \SI{44}{\micro\farad} nominal total sits inside the \SIrange{22}{68}{\micro\farad}
recommendation (\cref{sec:ap632xx}). Assumption: the parts are X5R/X7R ceramics rated
$\geq \SI{10}{\volt}$, retaining $\geq 60\,\%$ of nominal capacitance at 5~V DC bias ---
effective $C_{\mathrm{out}} \geq \SI{26}{\micro\farad}$, still above the recommendation's
floor. The datasheet's load-transient sizing bound (\ds{AP632XX.pdf}{14}), with a
permitted overshoot $\Delta V = \SI{50}{\milli\volt}$ and the worst CCM-formula peak:
\begin{align*}
  C_{\mathrm{out}} &> \frac{L\left(I_{\mathrm{out}}+\Delta I_L/2\right)^2}
                          {(\Delta V + V_{OUT})^2 - V_{OUT}^2}
    = \frac{\SI{4.7}{\micro\henry} \times (0.100+0.298)^2\,\si{\square\ampere}}{5.05^2 - 5^2\ \si{\square\volt}}
    = \frac{\SI{0.744}{\micro\joule}}{\SI{0.503}{\square\volt}}
    = \SI{1.5}{\micro\farad},
\end{align*}
met with $17\times$ margin even after derating --- unsurprising, since the load has no
transient profile to speak of. Output ripple in the occupied PFM mode is bounded by one
pulse's full charge dumped into the (derated) capacitance plus the ESR step:
\begin{align*}
  \Delta V_{\mathrm{ripple}} &\lesssim \frac{Q_{\max}}{C_{\mathrm{out}}}
      + I_{\mathrm{pk}}\,\mathrm{ESR}
    = \frac{\SI{231}{\nano\coulomb}}{\SI{26}{\micro\farad}}
      + 0.45 \times \SI{2.5}{\milli\ohm}
    \approx 8.9 + 1.1 \approx \SI{10}{\milli\volt}\ \text{pk-pk},
\end{align*}
(assumed effective ESR \SI{2.5}{\milli\ohm} for two paralleled ceramics; the
fixed-frequency ripple, $\Delta I_L/(8 f_{SW} C) + \Delta I_L\,\mathrm{ESR} \approx
\SI{4}{\milli\volt}$, is smaller still). Stacking the worst case against the load's
requirement: the AP63205 regulates to \SIrange{4.95}{5.05}{\volt} over conditions
(\cref{sec:ap632xx}), and $4.95 - 0.010 = \SI{4.94}{\volt}$ /
$5.05 + 0.010 = \SI{5.06}{\volt}$ sit inside the K155ID1's recommended
\SIrange{4.75}{5.25}{\volt} supply window (\cref{sec:bcd-driver}) with
$\approx\SI{190}{\milli\volt}$ margin on both sides. Match.

\subsection{Input capacitance}\label{sec:5v-cin}

The populated C401 (\SI{10}{\micro\farad} ceramic) equals the datasheet's BOM-table
example; the accompanying prose asks for ``greater than \SI{10}{\micro\farad}''
(\cref{sec:ap632xx}), so the population meets the manufacturer's own example rather than
strictly exceeding the prose figure --- noted for precision, not flagged, because the
requirement scales with load and this rail runs at 1/20th of the part's rating. The
buck's input RMS ripple current at the worst point ($D$ nearest 0.5, i.e.\
$V_{IN}=\SI{8.5}{\volt}$):
\begin{align*}
  I_{\mathrm{RMS}} &\approx I_{\mathrm{out}}\sqrt{D(1-D)}
    = 0.100\sqrt{0.588 \times 0.412} = \SI{49}{\milli\ampere},
\end{align*}
against the guidance that the capacitor's RMS rating exceed half the maximum load
(\SI{50}{\milli\ampere} for this application; ampere-class for any ceramic
\SI{10}{\micro\farad}) --- margin of order $50\times$. Input ripple voltage at nominal
capacitance is $I_{\mathrm{out}} D(1-D)/(f_{SW} C) \approx \SI{2}{\milli\volt}$, still
$\approx\SI{3}{\milli\volt}$ after DC-bias derating at 13~V. The one thing the schematic
does not record is the capacitor's voltage rating: on a 13~V rail it must be
$\geq\SI{16}{\volt}$, and \SI{25}{\volt} is preferred for bias-derating headroom ---
confirm at BOM finalization, same status as the inductor rating above. C402
(\SI{100}{\nano\farad}, SW-to-BST) matches the recommended bootstrap value exactly.
Match throughout.

\subsection{Soft start, enable, and dissipation}\label{sec:5v-ss}

Soft start is the fixed internal \SI{4}{\milli\second} ramp, not externally programmable
(\cref{sec:ap632xx}); there is no component to verify. Start-up charging current is
correspondingly trivial:
\begin{align*}
  I_{\mathrm{chg}} &= C_{\mathrm{out}}\frac{dV}{dt}
    = \SI{44}{\micro\farad} \times \frac{\SI{5}{\volt}}{\SI{4}{\milli\second}}
    = \SI{55}{\milli\ampere},
\end{align*}
which, even summed with the full \SI{100}{\milli\ampere} load, stays a factor of 16 below
the minimum current limit --- no inrush interaction with the limit or with hiccup mode.
EN tied directly to $V_{\mathrm{BUS}}$ is one of the two sanctioned always-on
configurations (\cref{sec:ap632xx}): the converter starts as soon as the upstream load
switch presents $V_{\mathrm{BUS}}$, with the drivers' 5~V arriving about
\SI{4}{\milli\second} later, and no sequencing constraint exists among this rail, the
3.3~V rail, or the flyback at this level. Cross-rail power-up ordering, to the extent it
matters at all, is a system-level topic taken up after the individual rails are
established, not a property of this converter.

Dissipation, using the conservative $\eta = 0.85$ from \cref{sec:5v-op}:
\begin{align*}
  P_{\mathrm{diss}} &= P_{\mathrm{out,max}}\left(\frac{1}{\eta}-1\right)
    = 0.50 \times 0.176 = \SI{88}{\milli\watt}, &
  \Delta T &= \theta_{JA} P_{\mathrm{diss}} = 89 \times 0.088 \approx \SI{8}{\kelvin},
\end{align*}
so $T_J \approx \SI{48}{\celsius}$ at a \SI{40}{\celsius} ambient --- far below the
\SI{125}{\celsius} design target adopted in \cref{sec:ap632xx}. No thermal management
beyond the footprint's own copper is required.

\subsection{Summary}\label{sec:5v-summary}

Assumptions carried through this section: inductor $I_{\mathrm{sat}} \geq \SI{3.0}{\ampere}$
and DCR $\leq \SI{50}{\milli\ohm}$ (generic 6045-class part, no part number recorded);
output ceramics retain $\geq 60\,\%$ capacitance at 5~V bias with effective ESR
\SI{2.5}{\milli\ohm}; converter efficiency 0.85 (conservative floor under the datasheet's
$\approx$0.88--0.90 typical curve); indicator-LED forward voltage \SI{2.0}{\volt};
PFM pulses idealized as clean
\SI{450}{\milli\ampere}-peak triangles for the ripple and pulse-rate estimates. Each
appears at its point of use above.

\begin{longtable}{@{}p{0.26\textwidth}p{0.30\textwidth}p{0.22\textwidth}p{0.14\textwidth}@{}}
\toprule
Parameter & Derived requirement & Populated & Verdict \\
\midrule
\endfirsthead
\toprule
Parameter & Derived requirement & Populated & Verdict \\
\midrule
\endhead
\bottomrule
\endlastfoot
Load current & $4\times I_{CC}$ (64/100~mA) $+$ D401 indicator
  (\SI{0.59}{\milli\ampere} via R401): \SI{65}{\milli\ampere} typ /
  \SI{101}{\milli\ampere} max, static & --- & Derived \\
Input window & inside \SIrange{3.8}{32}{\volt} recommended &
  $V_{\mathrm{BUS}}$ \SIrange{8.5}{13}{\volt}, direct & Match (large margin) \\
Duty cycle & 0.385--0.588; internal slope compensation covers $D>0.5$ & --- &
  Verified \\
Conduction mode & load $<$ CCM boundary (\SIrange{199}{298}{\milli\ampere})
  $\rightarrow$ automatic PFM; regulation maintained & --- & Verified (benign) \\
Minimum on-time & \SI{264}{\nano\second} (PFM) / \SI{330}{\nano\second} (PWM,
  $+6\%$ dither) vs.\ \SI{80}{\nano\second} & --- & Match ($\geq 3.3\times$) \\
Inductance L401 & Eq.-7 window \SIrange{2.8}{4.7}{\micro\henry} at \SI{13}{\volt};
  datasheet 5~V BOM value \SI{4.7}{\micro\henry} & \SI{4.7}{\micro\henry}
  & Match (reference BOM followed) \\
Inductor $I_{\mathrm{sat}}$ & $\geq\SI{0.45}{\ampere}$ normal;
  $\geq\SI{3.1}{\ampere}$ covers fault peak & 6045-class, rating unrecorded
  (assumed $\geq\SI{3}{\ampere}$) & Match under assumption; \textbf{confirm part} \\
Current-limit margin & peak $\leq\SI{0.45}{\ampere}$ vs.\ \SI{2.5}{\ampere} min
  limit & --- & Large margin ($5.6\times$) \\
Output capacitance C403/C404 & \SIrange{22}{68}{\micro\farad}
  ($\geq\SI{1.5}{\micro\farad}$ transient bound) & $2\times\SI{22}{\micro\farad}$
  (= BOM example) & Match \\
Output ripple & $\lesssim\SI{10}{\milli\volt}$ pk-pk (PFM worst case) & --- &
  Derived; $\ll$ load tolerance \\
Rail accuracy at load & K155ID1 requires \SIrange{4.75}{5.25}{\volt} &
  \SIrange{4.94}{5.06}{\volt} incl.\ ripple & Match ($\approx\SI{190}{\milli\volt}$
  both sides) \\
Input capacitor C401 & $\geq\SI{10}{\micro\farad}$ ceramic; RMS
  $\geq\SI{49}{\milli\ampere}$ & \SI{10}{\micro\farad} (= BOM example; rating
  unrecorded) & Match; \textbf{confirm $\geq$16~V rating} \\
BST capacitor C402 & \SI{100}{\nano\farad} SW-to-BST & \SI{100}{\nano\farad} &
  Match \\
Soft start & \SI{4}{\milli\second} fixed internal; inrush
  $\leq\SI{155}{\milli\ampere}$ & internal, no component & Match \\
EN configuration & EN\,=\,VIN valid always-on & EN tied to $V_{\mathrm{BUS}}$ &
  Match \\
Converter dissipation & \SI{88}{\milli\watt}, $\Delta T \approx \SI{8}{\kelvin}$
  vs.\ \SI{125}{\celsius} target & TSOT26, $\theta_{JA}=89$~\si{\celsius\per\watt} &
  Match (large margin) \\
Bypass C1--C4 & \SI{100}{\nano\farad} per TTL package, standard practice &
  \SI{100}{\nano\farad} $\times 4$ & Match \\
\end{longtable}

No populated value is flagged as a mismatch. The two open items are both unrecorded
ratings rather than wrong values: the actual L401 part's saturation current and the C401
voltage rating ($\geq\SI{16}{\volt}$ on a 13~V rail), both to be confirmed when the BOM
is finalized. The rail otherwise reduces to the manufacturer's reference design operating
at one-twentieth of its rating into a static load --- every margin is wide, and the
section's only substantive finding is the (benign, automatic) PFM operating mode.
